% GENERATED FILE. Edit canonical lessons, then run npm run generate:books.
% canonical-source-hash: fnv1a64:1f3c61ce382043d2
% canonical-lessons: RU-C232-kopit, RU-C232-zanimat, RU-C232-zvonit, RU-C232-otpravlyat, RU-C232-poluchat

\chapter{To save up, To borrow, To ring, To send, To receive}
\label{ch:ru-to-save-up-to-borrow-to-ring-to-send-to-receive}
\hlchaptermodality{232}

\begin{chapteropening}
\textbf{By the end of this chapter:} \emph{I can say to save up, to borrow, to ring, to send and to receive.}

Complete the last lesson of chapter 232: I can say to save up, to borrow, to ring, to send and to receive.
\end{chapteropening}

\section[\texorpdfstring{\ru{копить} (kopít)}{копить (kopít)}]{\ru{копить} (kopít) — to save up}

\label{lesson:RU-C232-kopit}



\begin{quote}
Before the new one: say the Russian for a class, a lesson, then the Russian for a mark, a grade.
\end{quote}

\subsection*{You'll want to know: \ru{копить}}
\textbf{\ru{копить}} (\emph{kopít}) — \textquotedblleft{}to save up\textquotedblright{}.

\begin{grammarlens}[title={What you've built}]
One more word for this chapter.
\end{grammarlens}

\subsection*{Your turn}
\begin{itemize}
\raggedright
  \item \emph{Say it:} \emph{kopít}
  \item \emph{Say it:} \emph{kopít}, once more
  \item \emph{Say it:} say \emph{kopít}
  \item \emph{Recall:} say the Russian for a class, a lesson, then the Russian for a mark, a grade, then say \emph{kopít} again
\end{itemize}

\begin{tcolorbox}[breakable,colback=teal!4,colframe=teal!35!black,title={Before you move on}]
What does \ru{копить} mean? (\textquotedblleft{}To save up\textquotedblright{}.) Say it once more.
\end{tcolorbox}
\section[\texorpdfstring{\ru{занимать} (zanimát)}{занимать (zanimát)}]{\ru{занимать} (zanimát) — to borrow; to take up (space)}

\label{lesson:RU-C232-zanimat}



\begin{quote}
Before the new one: say the Russian for a mark, a grade, then the Russian for to save up.
\end{quote}

\subsection*{You'll want to know: \ru{занимать}}
\textbf{\ru{занимать}} (\emph{zanimát}) — \textquotedblleft{}to borrow; to take up (space)\textquotedblright{}.

\begin{grammarlens}[title={What you've built}]
One more word for this chapter.
\end{grammarlens}

\subsection*{Your turn}
\begin{itemize}
\raggedright
  \item \emph{Say it:} \emph{zanimát}
  \item \emph{Say it:} \emph{zanimát}, once more
  \item \emph{Say it:} say \emph{zanimát}
  \item \emph{Recall:} say the Russian for a mark, a grade, then the Russian for to save up, then say \emph{zanimát} again
\end{itemize}

\begin{tcolorbox}[breakable,colback=teal!4,colframe=teal!35!black,title={Before you move on}]
What does \ru{занимать} mean? (\textquotedblleft{}To borrow; to take up (space)\textquotedblright{}.) Say it once more.
\end{tcolorbox}
\section[\texorpdfstring{\ru{звонить} (zvonít)}{звонить (zvonít)}]{\ru{звонить} (zvonít) — to ring, to phone}

\label{lesson:RU-C232-zvonit}



\begin{quote}
Before the new one: say the Russian for to save up, then the Russian for to borrow; to take up.
\end{quote}

\subsection*{You'll want to know: \ru{звонить}}
\textbf{\ru{звонить}} (\emph{zvonít}) — \textquotedblleft{}to ring, to phone\textquotedblright{}.

\begin{grammarlens}[title={What you've built}]
One more word for this chapter.
\end{grammarlens}

\subsection*{Your turn}
\begin{itemize}
\raggedright
  \item \emph{Say it:} \emph{zvonít}
  \item \emph{Say it:} \emph{zvonít}, once more
  \item \emph{Say it:} say \emph{zvonít}
  \item \emph{Recall:} say the Russian for to save up, then the Russian for to borrow; to take up, then say \emph{zvonít} again
\end{itemize}

\begin{tcolorbox}[breakable,colback=teal!4,colframe=teal!35!black,title={Before you move on}]
What does \ru{звонить} mean? (\textquotedblleft{}To ring, to phone\textquotedblright{}.) Say it once more.
\end{tcolorbox}
\section[\texorpdfstring{\ru{отправлять} (otpravlyát)}{отправлять (otpravlyát)}]{\ru{отправлять} (otpravlyát) — to send}

\label{lesson:RU-C232-otpravlyat}



\begin{quote}
Before the new one: say the Russian for to borrow; to take up, then the Russian for to ring.
\end{quote}

\subsection*{You'll want to know: \ru{отправлять}}
\textbf{\ru{отправлять}} (\emph{otpravlyát}) — \textquotedblleft{}to send\textquotedblright{}.

\begin{grammarlens}[title={What you've built}]
One more word for this chapter.
\end{grammarlens}

\subsection*{Your turn}
\begin{itemize}
\raggedright
  \item \emph{Say it:} \emph{otpravlyát}
  \item \emph{Say it:} \emph{otpravlyát}, once more
  \item \emph{Say it:} say \emph{otpravlyát}
  \item \emph{Recall:} say the Russian for to borrow; to take up, then the Russian for to ring, then say \emph{otpravlyát} again
\end{itemize}

\begin{tcolorbox}[breakable,colback=teal!4,colframe=teal!35!black,title={Before you move on}]
What does \ru{отправлять} mean? (\textquotedblleft{}To send\textquotedblright{}.) Say it once more.
\end{tcolorbox}
\section[\texorpdfstring{\ru{получать} (poluchát)}{получать (poluchát)}]{\ru{получать} (poluchát) — to receive, to get}

\label{lesson:RU-C232-poluchat}



\begin{quote}
Before the new one: say the Russian for to ring, then the Russian for to send.
\end{quote}

\subsection*{You'll want to know: \ru{получать}}
\textbf{\ru{получать}} (\emph{poluchát}) — \textquotedblleft{}to receive, to get\textquotedblright{}.

\begin{grammarlens}[title={What you've built}]
One more word for this chapter.
\end{grammarlens}

\subsection*{Your turn}
\begin{itemize}
\raggedright
  \item \emph{Say it:} \emph{poluchát}
  \item \emph{Say it:} \emph{poluchát}, once more
  \item \emph{Say it:} say \emph{poluchát}
  \item \emph{Recall:} say the Russian for to ring, then the Russian for to send, then say \emph{poluchát} again
\end{itemize}

\begin{tcolorbox}[breakable,colback=teal!4,colframe=teal!35!black,title={Before you move on}]
What does \ru{получать} mean? (\textquotedblleft{}To receive, to get\textquotedblright{}.) Say it once more.
\end{tcolorbox}
